The goal of the following exercises is to (almost) give a proof of the fact that
anodyne maps are precisely those monomorphisms which are weak equivalences.
More precisely, we will show that it is implied by the following statement: Let
$p\colon X\to Y$ be a Kan fibration. Then there exists a factorization of $p$ as
\[
 X\xrightarrow{\alpha}Z\xrightarrow{\beta}Y
\]
where $\beta$ is a trivial fibration (i.e., it has the RLP with respect to monomorphisms)
and $\alpha$ is a minimal fibration. In this context, we need to know the following facts
about minimal fibrations:
\begin{itemize}
 \item A minimal fibration is a Kan fibration.
 \item A minimal fibration $\alpha\colon X\to Z$ is locally trivial, i.e., for every simplex
       $\Delta^n\to Z$, the pulled-back fibration is isomorphic (over $\Delta^n$) to a
       projection $\Delta^n\times B\to\Delta^n$.
\end{itemize}

\begin{exercise}{107}
Show that a Kan fibration $p\colon X\to Y$ is a trivial fibration if and only if its
fibres are contractible. Hint: For the interesting direction, consider a lifting problem
\[
\begin{tikzcd}[column sep=large,row sep=large]
 \partial\Delta^n\arrow[r,"a"]\arrow[d] & X\arrow[d] \\
 \Delta^n\arrow[r,"b"']\arrow[ur,dashed] & Y
\end{tikzcd}
\]
For future reference, we let $y=b(n)$. Consider the map
$g\colon\Delta^n\times\Delta^1\to\Delta^n$ determined by
$g(k,0)=k$ and $g(k,1)=n$. Show that the map
$\partial\Delta^n\times\{0\}\to\partial\Delta^n\times\Delta^1$ is anodyne, and
use $g$ to obtain a map
\[
 \Delta^1\to X^{\partial\Delta^n}\times_{Y^{\partial\Delta^n}}Y^{\Delta^n}
\]
sending $0$ to the original square and $1$ to the square
\[
\begin{tikzcd}[column sep=large,row sep=large]
 \partial\Delta^n\arrow[rr]\arrow[d] & & X\arrow[d] \\
 \Delta^n\arrow[r]\arrow[urr,dashed] & \Delta^0\arrow[r,"y"] & Y
\end{tikzcd}
\]
where the map $\Delta^0\to Y$ is the object $y$. Show that this lifting problem can
be solved. Finally, show that this implies that the original lifting problem can be
solved as well.

\begin{proof}
If $p$ is a trivial fibration, each fibre $F\to\Delta^0$ is again a
trivial fibration. Since $F$ is Kan, Exercise~\ref{ex:63} shows that
$F$ is contractible.

For the converse, we use two consequences of Lemma~1.3.31(4).
For every simplicial set $K$, the map
$K\times\{0\}\hookrightarrow K\times\Delta^1$ is anodyne.
Also, for $n\ge0$, the inclusion
\[
 C_n=(\partial\Delta^n\times\Delta^1)\cup(\Delta^n\times\{1\})
 \hookrightarrow\Delta^n\times\Delta^1
\]
is anodyne. These are the pushout products of
$\varnothing\hookrightarrow K$ with $\{0\}\hookrightarrow\Delta^1$,
and of $\partial\Delta^n\hookrightarrow\Delta^n$ with
$\{1\}\hookrightarrow\Delta^1$, respectively.
The endpoint inclusions are horn inclusions, so the lemma applies.

First, a contractible Kan complex $F$ fills every boundary
$u\colon\partial\Delta^n\to F$.
Indeed, choose a contraction $h\colon F\times\Delta^1\to F$ from
the identity to a constant vertex.
On $C_n$, use $h(u(-),-)$ on the boundary cylinder and this constant
vertex on $\Delta^n\times\{1\}$.
The maps agree on their intersection. The Kan condition extends the
resulting map across the anodyne inclusion $C_n\hookrightarrow
\Delta^n\times\Delta^1$; its restriction at $0$ fills $u$.
This includes $n=0$, when contractibility supplies a vertex.

Now take the square $(a,b)$ in the question and use the specified map
$g$ and vertex $y=b(n)$.
The anodyne inclusion
$\partial\Delta^n\times\{0\}\hookrightarrow
\partial\Delta^n\times\Delta^1$ lifts against $p$, giving
$H_\partial\colon\partial\Delta^n\times\Delta^1\to X$ with
$H_\partial(-,0)=a$ and $pH_\partial=bg$ on this cylinder.
The adjoints of $H_\partial$ and $bg$ give exactly the path of squares
requested in the hint.
At $1$, the boundary lies in the fibre $X_y$.
By the preceding paragraph it extends to a map
$c\colon\Delta^n\to X_y$.
Glue $H_\partial$ to $c$ on $\Delta^n\times\{1\}$.
This defines a map $C_n\to X$ over $bg$, and another application of
the Kan lifting property gives
$L\colon\Delta^n\times\Delta^1\to X$ over $bg$ extending it.
Then $L(-,0)$ extends $a$ and lies over $b$, as required.
Thus $p$ lifts every boundary inclusion and is a trivial fibration.
If $Y=\varnothing$, then also $X=\varnothing$ and the conclusion is immediate.
\end{proof}
\end{exercise}
